\section{Methodology}
\label{sec:method}

We adapt the \skillmix paradigm to creative story generation through a three-stage pipeline: component extraction, combinatorial specification, and controlled generation.

\subsection{Component Extraction}
\label{sec:extraction}

We sample 100 stories from the \rocstories training set~\cite{mostafazadeh2016corpus} (78,528 five-sentence commonsense stories, seed=42) and extract structured components using \gptfour with temperature 0.2 and structured JSON output.
For each story, we extract six components:

\begin{itemize}[leftmargin=*,itemsep=0pt,topsep=0pt]
    \item \textbf{Setting}: location and temporal context
    \item \textbf{Character}: protagonist description and role
    \item \textbf{Conflict type}: categorical (person vs.\ self, nature, person, society, fate, or technology)
    \item \textbf{Theme}: central thematic concern
    \item \textbf{Emotional arc}: trajectory of emotional tone (positive resolution, negative resolution, bittersweet, neutral, comedic)
    \item \textbf{Tone}: overall narrative tone (heartwarming, serious, lighthearted, melancholic, dramatic, suspenseful)
\end{itemize}

All 100 stories were successfully decomposed with valid, well-formed JSON output (100\% extraction rate).

\subsection{Component Taxonomy}
\label{sec:taxonomy}

The extracted components reveal meaningful distributional patterns (\tabref{tab:taxonomy}).
Person-vs-self conflicts dominate (54\%), reflecting \rocstories' focus on everyday personal dilemmas.
Positive resolution is the most common arc (53\%), and heartwarming the most frequent tone (35\%).
One limitation is that narrative technique was classified as linear for all 100 stories---\rocstories' five-sentence format is too short to exhibit flashbacks, frame narratives, or in medias res structures.

\subsection{Combinatorial Specification}
\label{sec:combination}

To create novel component specifications, we sample each of the six components from a \emph{different} source story, producing combinations that are unlikely to occur naturally.
For each of the 50 component-controlled stories, we sample five source stories and draw one component from each (with the sixth drawn from a random additional source).
This is analogous to the skill-mixing step in \skillmix, where novel combinations of individually observed skills create training examples that test compositional generalization.

\subsection{Story Generation}
\label{sec:generation}

We generate stories under three conditions using \gptfour with temperature 0.8:

\para{\compcontrolled ($n{=}50$).}
The model receives a full 7-component specification (setting, character, conflict, theme, arc, tone, plus a title derived from the theme) and must produce a story satisfying all constraints.

\para{\titleonly baseline ($n{=}50$).}
The model receives only the theme as a title and generates a story freely.

\para{\prompted baseline ($n{=}50$).}
The model receives the theme and setting, providing minimal structural guidance.

Additionally, we run a supplementary experiment comparing \emph{coherent} specifications (all components from the same source story, $n{=}25$) against \emph{mixed} specifications (components from different sources, $n{=}25$) to isolate the effect of component incoherence.

\subsection{Evaluation}
\label{sec:evaluation}

\para{Quality assessment.}
We evaluate all stories using \gptfour as judge (temperature 0.1) on five dimensions, each scored 1--5: coherence, creativity, character quality, language quality, and overall quality.

\para{Controllability assessment.}
For component-controlled stories, we additionally score adherence to each specified component (setting, character, conflict, theme, arc, tone) and compute an overall adherence score, each on a 1--5 scale.

\para{Diversity metrics.}
We compute three automatic diversity measures across each condition's story set: unique trigram ratio, vocabulary size, and mean pairwise inter-story trigram overlap.

\para{Statistical analysis.}
We use Wilcoxon signed-rank tests (paired, non-parametric) with Bonferroni correction for quality comparisons between conditions. We report Cohen's $d$ effect sizes for all pairwise comparisons. Significance threshold is $\alpha = 0.05$ after correction.
